\documentclass[10pt]{article}
\usepackage[a4paper,margin=25mm]{geometry}
\usepackage{amsmath,amssymb,amsthm}
\usepackage[hidelinks]{hyperref}
\newtheorem{theorem}{Theorem}
\setlength{\emergencystretch}{1em}
\hypersetup{hidelinks}
\begin{document}
\title{LYM posets are not closed under arbitrary direct products}
\author{Chengzhe Gao}\date{}\maketitle
For a finite ranked poset $P$ with rank levels $P_i$, the LYM property is
\[
 \sum_{x\in A}\frac1{|P_{\rho(x)}|}\leq1
 \quad\text{for every antichain }A\subseteq P.
\]
Conjecture 2594 asserts that the direct product of two LYM posets always
has this property, without additional hypotheses on their rank sizes.

Let $P$ have elements $a,b,c,d,e$ with ranks $0,1,2,2,2$. Put every element
of a lower rank below every element of a higher rank, and leave distinct
elements of rank two incomparable. Thus $P$ is the ordinal sum of antichains
of sizes $(1,1,3)$. Let $Q=\{0<1\}$ be a two-element chain.

\begin{theorem}
Both $P$ and $Q$ are LYM posets, but $P\times Q$ is not.
\end{theorem}
\begin{proof}
An antichain of $P$ lies within a single rank, because any two elements
of different ranks are comparable. Its Lubell sum is therefore its
cardinality divided by the size of that rank, at most one. The same
argument applies to $Q$.

Use componentwise order and additive rank on $P\times Q$. Its rank sizes
are the coefficients of
\[
 (1+t+3t^2)(1+t)=1+2t+4t^2+3t^3.
\]
Consider
\[
 A=\{(a,1),(c,0),(d,0),(e,0)\}.
\]
The last three elements are mutually incomparable. Each is incomparable
with $(a,1)$, since its first coordinate is greater and its second is
smaller. The rank of $(a,1)$ is one, and the ranks of the other three
elements are two. Consequently its Lubell sum is
\[
 \frac12+\frac34=\frac54>1.
\]
This violates LYM in the actual direct product.
\end{proof}

The antichain is also maximal: the omitted elements $(a,0),(b,0)$ lie
below one of its members, while $(b,1),(c,1),(d,1),(e,1)$ lie above one.
Thus wording in the source referring to maximal antichains does not affect
the counterexample. The further claims about matching permutations are
unnecessary once the asserted product closure fails.

\paragraph{Formal verification.}
\texttt{RankedPoset} includes the full enumerated vertex universe, its
no-duplication and completeness proofs, the three partial-order axioms,
rank zero on minimal elements, the common maximal rank, and rank increment
one on every covering pair. All these requirements are verified for
$P,Q,R$, where $R$ uses actual ordered pairs and the componentwise relation.
\texttt{product\_\allowbreak{}order} and \texttt{product\_\allowbreak{}rank} explicitly identify
$R$ with the standard ranked direct product. Rank sizes are counted from
the complete vertex lists.

\texttt{lubell} uses \texttt{Std.Internal.Rat}, performing the ordinary
rational sum of reciprocal rank sizes. The two LYM proofs cover every
Boolean membership function, by checking all $32$ or $4$ tuples and proving
the respective tuple expansion identities. The theorem
\texttt{every\_\allowbreak{}subset\_\allowbreak{}represented} represents every arbitrary
propositional subset by its characteristic function, and
\texttt{lym\_\allowbreak{}for\_\allowbreak{}every\_\allowbreak{}subset} gives the corresponding LYM statement.
Thus the enumeration omits no antichains. The explicit product witness,
its antichain property, its exact rational weight, and failure of the
LYM bound are all checked. The theorem
\texttt{conjecture2594\_\allowbreak{}false} negates the necessary instance
$\mathrm{LYM}(P)\Rightarrow\mathrm{LYM}(Q)\Rightarrow\mathrm{LYM}(R)$.

Lean 4.19.0 checks all computation by kernel reduction, without
\texttt{native\_\allowbreak{}decide}, unproved declarations, or additional axioms.

\end{document}
